% Laget automatisk av elevark/lag_elevark.py fra opplegg/eggslipp.html. Endringer her blir overskrevet.
\documentclass[11pt]{article}
\usepackage{elevark}
\begin{document}
\elevarkhode{Fysikk 1 · Bevegelsesmengde og støt}{Eggslippkonkurranse}{https://lektorbo.no/opplegg/eggslipp.html}{Opplegget på nett}

Lag en innpakning som gjør at et rått egg tåler å slippes fra så høyt som mulig, med så lite materiale som mulig. Fysikken bak er impuls: $ΣF \cdot Δt = Δp$.

\section*{Regler}

\begin{enumerate}
\item Hver gruppe får ett rått egg, 10 sugerør, 2 A4-ark, 1~m teip og 1~m hyssing. Andre materialer er ikke lov. Saks er lov.
\item Innpakningen skal lages på 30 minutter. Egget må kunne tas ut etterpå, så det kan sjekkes.
\item Innpakningen veies uten egg før første slipp.
\item Egget slippes fra 1,0~m, 2,0~m og 3,0~m, eller fra trappeavsatser med målt høyde. Gruppen velger selv når de vil stoppe.
\item Det er bare den høyeste høyden egget tålte som teller. Knuser egget, gjelder høyden før.
\end{enumerate}

\section*{Poeng}

$\displaystyle \text{poeng} = 10 \cdot \frac{\text{høyde i cm}}{\text{masse i gram}}$

En innpakning på 40~g som tåler 3,0~m, får $10 \cdot 300 / 40 = 75$ poeng. Lette innpakninger lønner seg, men bare hvis egget overlever.

\section*{Fysikken}

Egget har samme $Δp = mv$ uansett hva det lander på. Innpakningen kan bare gjøre én ting: forlenge bremsetiden $Δt$, slik at kraften blir mindre. Skriv inn høyden under for å se hvor lang bremsetid og bremsestrekning som trengs.

\section*{Etterarbeid}

\begin{oppgaveliste}
\oppgave \fase{test}{Regn}Regn ut farten egget hadde rett før landing fra den høyeste høyden dere klarte, og endringen i bevegelsesmengde.
\svar
\oppgave \fase{test}{Mål}Film landingen med sakte film. Hvor langt ble innpakningen presset sammen? Bruk det til å anslå bremsetiden og middelkraften på egget.
\svar
\oppgave \fase{forklar}{Forklar}Hva gjorde de beste innpakningene som de dårligste ikke gjorde? Bruk ordene impuls, bremsetid og kraft.
\svar
\oppgave \fase{forklar}{Forklar}Hvorfor hjelper det å ha egget i midten med mye «knasesone» rundt, og ikke bare et tykt lag på den ene siden?
\svar
\oppgave \fase{utfordring}{Utfordring}Luftmotstand kan også hjelpe. Hvor stor må fallskjermen være for at egget skal lande med 3~m/s? Bruk $L = \tfrac{1}{2}ρC_{\mathrm{d}}Av^{2}$ med $C_{\mathrm{d}} \approx 1{,}3$ og $ρ = 1{,}2\,\mathrm{kg/m^{3}}$.
\svar
\end{oppgaveliste}

\end{document}
